\documentclass[
  spanish,
  letterpaper, oneside
]{article}

\def\documenttitle {Ser estudiante en entornos virtuales de enseñanza-aprendizaje}
\def\documentsubtitle {Brújula de Interaprendizaje y compromisos de mejora}
\def\documentsubject {Actividad 4 - Interaprendizaje en ambientes virtuales}

\def\documentauthor {Martin Jonathan de la Cruz}
\def\coursename {Interaprendizaje en ambientes virtuales}
\def\coursecode {004}

\def\universityname {Universidad Abierta y a Distancia de México}
\def\universityfaculty {Licenciatura en Derecho}
\def\universitydepartment {Interaprendizaje en ambientes virtuales}
\def\universitydepartmentimage {departamentos/UnADM}
\def\universitydepartmentimagecfg {height=1.57cm}
\def\universitylocation {Roma Norte, Ciudad de México}

\def\authortable {
  \begin{tabular}{ll}
    \textbf{Alumno:}
    & \begin{tabular}[t]{l}
      Martin Jonathan de la Cruz \\
    \end{tabular} \\
    Matrícula: & ES2611202040 \\
    Figura docente: & Janeth Santana Ramírez \\
    Semestre/Bloque: & 2 / 1 \\
    & \\
    \multicolumn{2}{l}{Fecha de realización: \today} \\
    \multicolumn{2}{l}{\universitylocation}
  \end{tabular}
}

\input{template}
\usepackage{helvet}
\renewcommand{\familydefault}{\sfdefault}
\usepackage{array}
\usepackage{longtable}
\usepackage{booktabs}
\usepackage{pdflscape}
\usepackage{tikz}
\usetikzlibrary{arrows.meta,positioning,calc,matrix,fit,shapes.geometric,shadows.blur}
\setcitestyle{authoryear,open={(},close={)}}

\def\coverwatermarkenabled {true}
\def\coverwatermarkimage {img/departamentos/UnADM.pdf}
\def\coverwatermarkopacity {0.16}
\def\coverwatermarkwidth {0.72\paperwidth}
\def\coverwatermarkxshift {0cm}
\def\coverwatermarkyshift {-0.35cm}

\newcommand{\insertcoverwatermark}{%
  \ifthenelse{\equal{\coverwatermarkenabled}{true}}{%
    \AddToShipoutPictureBG*{%
      \begin{tikzpicture}[remember picture,overlay]
        \node[
          opacity=\coverwatermarkopacity,
          inner sep=0pt,
          xshift=\coverwatermarkxshift,
          yshift=\coverwatermarkyshift
        ] at (current page.center) {%
          \includegraphics[width=\coverwatermarkwidth]{\coverwatermarkimage}%
        };
      \end{tikzpicture}%
    }%
  }{}%
}

\definecolor{unadmGreenDark}{HTML}{174A3A}
\definecolor{unadmGreen}{HTML}{5F8F3A}
\definecolor{unadmGold}{HTML}{B88A2A}
\definecolor{unadmPaper}{HTML}{F6F7F2}
\definecolor{unadmInk}{HTML}{1F2A24}
\definecolor{unadmRowAlt}{HTML}{EEF2E6}

% -----------------------------------------------------------------------------
% AulaTeX:ficha-producto v1 -- bloque de la Ficha de recuperación.
% La ficha es el entregable con que se materializa la técnica CUESTIONARIO
% (#32 de las 100 Técnicas Didácticas UnADM): el instrumento recoge datos y la
% ficha los convierte en lectura razonada y compromisos verificables.
% Caja institucional homóloga al wikibox de la materia.
% -----------------------------------------------------------------------------
\usepackage[most]{tcolorbox}
\usepackage{attachfile}
\usepackage{enumitem}

% attachfile: el 'color' debe darse como TRIPLE RGB numérico (0..1); un nombre
% \definecolor[HTML] provoca 'Argument of \atfi@textcolor has an extra }'.
\attachfilesetup{color={0.373 0.561 0.227},print=false}

% Botón de copia: appearance={} oculta el PushPin y muestra el 2.o argumento como
% botón visible. Adjunta el .txt embebido con la ficha copiable.
\newcommand{\fichaCopyButton}[1]{%
  \textattachfile[appearance={},description={Ficha de recuperacion (texto copiable)},mimetype=text/plain]{#1}{%
    \colorbox{unadmGreen}{\small\textcolor{white}{\,Copiar ficha\,}}%
  }%
}

\newtcolorbox{fichabox}[1][]{%
  enhanced, breakable,
  colback=unadmPaper, colframe=unadmGreen,
  boxrule=0pt, leftrule=3pt, arc=1.2mm,
  left=4mm, right=4mm, top=3mm, bottom=3mm,
  fonttitle=\bfseries\color{white}, coltitle=white, colbacktitle=unadmGreenDark,
  attach boxed title to top left={xshift=4mm,yshift=-2mm},
  boxed title style={colframe=unadmGreenDark,arc=0.6mm},
  #1
}

\newcommand{\fichameta}[4]{%
  \begin{tcolorbox}[colback=unadmRowAlt,colframe=unadmRowAlt,boxrule=0pt,arc=1mm,
    left=2mm,right=2mm,top=1.5mm,bottom=1.5mm]
  {\small\begin{tabular}{@{}p{3.4cm}p{9.6cm}@{}}
    \textbf{Instrumento} & #1 \\
    \textbf{Participante} & #2 \\
    \textbf{Fecha de aplicación} & #3 \\
    \textbf{Espacio / curso} & #4 \\
  \end{tabular}}
  \end{tcolorbox}\medskip}

\begin{document}

\insertcoverwatermark
\templatePortrait
\templatePagecfg
\onehalfspacing

\begin{abstractd}
  Ser estudiante en un entorno virtual no es una condición administrativa sino un
  conjunto de actitudes, habilidades y hábitos que se aprenden y pueden revisarse.
  Este documento ordena primero ese marco --autonomía y autogestión, comunicación
  escrita y dialógica, aprendizaje ubicuo mediado por TICCAD, colaboración e
  interaprendizaje, y metacognición atenta a la diversidad cultural y
  lingüística-- y después aplica la Brújula de Interaprendizaje, un cuestionario de
  autopercepción con escala de frecuencia que recorre esas cinco dimensiones. El
  registro razonado de las veinte afirmaciones permite identificar un perfil
  autogestivo pero poco recíproco, del que derivan dos compromisos verificables,
  uno personal y uno colaborativo, recogidos en la Ficha de recuperación.
\end{abstractd}

\templateIndex
\templateFinalcfg

\section{Introducción}

Nadie nace sabiendo estudiar en línea. La expresión ``ser estudiante en un entorno
virtual'' nombra algo más que estar inscrito en un aula digital: designa una forma de
aprender que se construye con el tiempo y que involucra decisiones sobre el propio
ritmo, la manera de comunicarse por escrito, el uso de las herramientas y la
disposición a aprender con otras personas. Por eso conviene tratarla como una
competencia revisable y no como un rasgo fijo.

El propósito de este trabajo es doble. Primero, precisar qué caracteriza al
estudiante de entornos virtuales según la literatura de la semana, con especial
atención a las dimensiones que sostienen el interaprendizaje. Después, aplicar esa
lectura a la propia experiencia mediante la Brújula de Interaprendizaje, un
instrumento de autopercepción cuyos resultados se recuperan en una ficha con
compromisos de mejora. El sentido de recorrer ambas partes es que el diagnóstico no
quede en la descripción: la autopercepción sirve si permite decidir qué cambiar y
cómo comprobar que cambió.

\section{Ser estudiante en entornos virtuales de enseñanza-aprendizaje}

\subsection{Una condición que se aprende}

El estudiante de entornos virtuales no es simplemente un estudiante presencial que
usa una computadora. La formación en línea reorganiza el tiempo, el espacio y los
canales de la interacción, de modo que exige actitudes y hábitos propios: iniciativa
para sostener el trabajo sin recordatorios continuos, disposición a comunicar por
escrito con precisión y capacidad para pedir ayuda antes de que el problema crezca
\citep{borges2007estudiante}. La didáctica universitaria en entornos virtuales
subraya que estas condiciones no dependen solo de la voluntad individual: se aprenden
cuando el diseño del curso las hace necesarias y visibles, y se debilitan cuando la
actividad se reduce a entregar productos \citep{bautista2016didactica}.

De ahí se sigue una consecuencia importante. Si la condición de estudiante en línea
se aprende, entonces puede diagnosticarse; y si puede diagnosticarse, puede mejorarse
mediante compromisos concretos. La autopercepción deja de ser una impresión difusa y
se convierte en el punto de partida de una decisión pedagógica.

\subsection{Las cinco dimensiones que se exploran}

La revisión de las lecturas permite ordenar en cinco dimensiones lo que se pone en
juego al estudiar en línea, y son precisamente las que recorre el instrumento
aplicado:

\begin{itemize}[leftmargin=1.6em,itemsep=0.35em]
  \item \textbf{Autogestión.} Organización del tiempo, constancia y capacidad de
  sostener el propio proceso sin depender de recordatorios externos.
  \item \textbf{Comunicación dialógica.} Calidad de la interacción escrita: no basta
  con participar, importa si la participación abre o cierra el intercambio.
  \item \textbf{Aprendizaje ubicuo.} Posibilidad de aprender en tiempos y lugares
  diversos gracias a las TICCAD, siempre que las herramientas se elijan por su función
  y no se acumulen \citep{zambranoPilay2020tecnologias}.
  \item \textbf{Colaboración.} Distinción entre el reparto de tareas y la construcción
  conjunta; el estándar de competencia sobre trabajo en equipo exige acuerdos
  explícitos, aportación efectiva y revisión del resultado colectivo, y no solo
  cumplimiento individual \citep{conocer2025ec0554}.
  \item \textbf{Metacognición atenta a la diversidad.} Revisión del propio aprendizaje
  articulada con el reconocimiento de la diversidad cultural y lingüística del grupo,
  dentro del marco de la educación intercultural \citep{aguadoSleeter2021educacion}.
\end{itemize}

Conviene precisar el alcance de la última dimensión, porque suele quedarse en el
enunciado. Adopto como criterio propio examinar si el lenguaje y los ejemplos que
utilizo facilitan efectivamente la participación de personas con trayectorias
distintas; es decir, la diversidad no se mide por la valoración que alguien declara,
sino por acciones concretas de ajuste del propio discurso.

\subsection{Por qué la autoevaluación es una práctica de interaprendizaje}

Podría objetarse que un instrumento de autopercepción es un ejercicio individual y,
por tanto, ajeno al interaprendizaje. La objeción se disuelve al observar qué evalúa:
cuatro de las cinco dimensiones describen relaciones con otras personas. Reconocer
los propios hábitos es la condición para saber qué se aporta y qué se obstaculiza en
la construcción conjunta. La experiencia de los interaprendizajes entre educadores
indígenas muestra justamente eso: el conocimiento comunitario se produce cuando cada
quien reconoce su lugar, lo pone en juego y lo transforma en la relación con las
demás personas \citep{bertely2011interaprendizajes}. La autoevaluación, en ese
sentido, no aísla: prepara para participar mejor.

\section{Aplicación de la Brújula de Interaprendizaje}

\subsection{El instrumento y su lógica}

La Brújula de Interaprendizaje es un cuestionario de autopercepción. En este
trabajo, su aplicación queda documentada en la ficha de recuperación y en la tabla de
registro: la primera identifica a la persona participante y el contexto, y la segunda
reúne las afirmaciones agrupadas por dimensión con la frecuencia marcada. Cada afirmación se responde con una escala de
frecuencia de cinco puntos --siempre, casi siempre, a veces, casi nunca, nunca--, es
decir, una escala tipo Likert que traduce hábitos en un dato ordenable y comparable
consigo mismo a lo largo del tiempo.

Esa estructura explica por qué el instrumento sirve para orientar y no para calificar.
La escala no produce un puntaje aprobatorio, sino un perfil: muestra en qué
dimensiones la conducta es sostenida y en cuáles es intermitente. Las veinte
afirmaciones se distribuyen en cuatro por dimensión, de manera que ninguna
dimensión pese más que otra y el perfil resultante sea legible como un conjunto de
tendencias, no como un promedio único.

\subsection{Registro razonado de las respuestas}

La tabla siguiente conserva la respuesta marcada en cada afirmación y añade la
lectura que de ella se desprende. Se incluye esa tercera columna porque el dato
aislado no dice nada: la frecuencia declarada solo adquiere sentido cuando se
interpreta contra el marco de la sección anterior. Así, marcar ``a veces'' en una
afirmación sobre revisión conjunta no significa un déficit de esfuerzo, sino que la
colaboración se está entendiendo como cumplimiento de la parte propia.

\clearpage
\pagestyle{empty}
\begin{landscape}
\thispagestyle{empty}
\begingroup
\fontsize{7.6}{8.6}\selectfont
\setlength{\tabcolsep}{3.2pt}
\renewcommand{\arraystretch}{1.12}
\begin{longtable}{@{}>{\centering\arraybackslash}p{0.04\linewidth}>{\raggedright\arraybackslash}p{0.40\linewidth}>{\centering\arraybackslash}p{0.13\linewidth}>{\raggedright\arraybackslash}p{0.38\linewidth}@{}}
\caption{Brújula de Interaprendizaje: registro de respuestas y lectura por dimensión}\label{tab:brujula-interaprendizaje-a4}\\
\toprule
\textbf{No.} & \textbf{Afirmación} & \textbf{Frecuencia marcada} & \textbf{Lectura del hallazgo} \\
\midrule
\endfirsthead
\toprule
\textbf{No.} & \textbf{Afirmación} & \textbf{Frecuencia marcada} & \textbf{Lectura del hallazgo} \\
\midrule
\endhead
\bottomrule
\endfoot
\bottomrule
\endlastfoot
% Escala de frecuencia (comentada, no visible): Siempre / Casi siempre / A veces / Casi nunca / Nunca.
\multicolumn{4}{@{}l}{\cellcolor{unadmRowAlt}\textbf{Dimensión 1. Autogestión y organización del estudio}} \\
\midrule
1 & Planifico mi semana de estudio a partir del calendario del aula antes de que inicie. & \textbf{Siempre} & La respuesta indica una práctica declarada de planeación anticipada; sin embargo, por sí sola no permite comparar esta conducta con las demás dimensiones ni inferir independencia respecto de recordatorios externos \citep{borges2007estudiante}. \\
\midrule
2 & Sostengo un ritmo constante de trabajo y no concentro el esfuerzo en el día de entrega. & \textbf{Casi siempre} & La constancia existe, pero cede cuando se acumulan entregas de varias asignaturas. \\
\midrule
3 & Dejo un margen de tiempo de reserva para imprevistos antes de la fecha de cierre. & \textbf{A veces} & Sin holgura, cualquier contratiempo se traslada íntegro a la calidad del producto. \\
\midrule
4 & Dedico tiempo a repasar lo aprendido aunque no haya una entrega inmediata. & \textbf{Casi nunca} & Lo primero que se sacrifica al apretar el calendario es la consolidación, no la entrega. \\
\midrule
% Escala de frecuencia (comentada, no visible): Siempre / Casi siempre / A veces / Casi nunca / Nunca.
\multicolumn{4}{@{}l}{\cellcolor{unadmRowAlt}\textbf{Dimensión 2. Comunicación e interacción dialógica}} \\
\midrule
5 & Participo en los espacios de interacción dentro de los plazos establecidos. & \textbf{Siempre} & La presencia está garantizada; la pregunta es de qué tipo es esa presencia. \\
\midrule
6 & Fundamento mis aportaciones con las fuentes revisadas en la asignatura. & \textbf{Siempre} & El aporte es sustantivo y verificable, condición de una comunicación académica útil. \\
\midrule
7 & Retomo de manera explícita las aportaciones de otras personas en mis mensajes. & \textbf{A veces} & Aquí se localiza la brecha central: la participación es expositiva más que conversacional \citep{bautista2016didactica}. \\
\midrule
8 & Formulo preguntas abiertas que inviten a otras personas a continuar el intercambio. & \textbf{Casi nunca} & Sin preguntas, el intercambio se agota en la suma de posiciones y no llega a diálogo. \\
\midrule
% Escala de frecuencia (comentada, no visible): Siempre / Casi siempre / A veces / Casi nunca / Nunca.
\multicolumn{4}{@{}l}{\cellcolor{unadmRowAlt}\textbf{Dimensión 3. Recursos digitales y aprendizaje ubicuo}} \\
\midrule
9 & Utilizo con soltura las herramientas del aula virtual y resuelvo por mi cuenta las fallas menores. & \textbf{Siempre} & El dominio técnico no es un obstáculo, de modo que las dificultades no son instrumentales. \\
\midrule
10 & Aprovecho tiempos y lugares distintos al escritorio para avanzar en el estudio. & \textbf{Casi siempre} & El aprendizaje ubicuo opera en la práctica y amplía la disponibilidad real de estudio \citep{zambranoPilay2020tecnologias}. \\
\midrule
11 & Elijo cada herramienta por la función que cumple y no por su disponibilidad o novedad. & \textbf{A veces} & La acumulación tecnológica supera a la integración: se guardan más recursos de los que se usan. \\
\midrule
12 & Integro en un solo lugar los materiales que reúno para poder recuperarlos después. & \textbf{Casi nunca} & La dispersión convierte el acopio en ruido y anula la ventaja de la ubicuidad. \\
\midrule
% Escala de frecuencia (comentada, no visible): Siempre / Casi siempre / A veces / Casi nunca / Nunca.
\multicolumn{4}{@{}l}{\cellcolor{unadmRowAlt}\textbf{Dimensión 4. Colaboración e interaprendizaje}} \\
\midrule
13 & Cumplo en tiempo y forma la parte del trabajo colectivo que me corresponde. & \textbf{Siempre} & La responsabilidad individual está cubierta, primer requisito del trabajo en equipo \citep{conocer2025ec0554}. \\
\midrule
14 & Participo en la definición de los acuerdos de trabajo del equipo antes de comenzar. & \textbf{Casi siempre} & Intervengo en el encuadre, lo que facilita que las reglas sean compartidas y no impuestas. \\
\midrule
15 & Reviso el producto completo del equipo, y no solo el apartado que redacté. & \textbf{A veces} & Sin revisión del conjunto, la entrega es una suma de partes y no una construcción conjunta. \\
\midrule
16 & Propongo un cierre colectivo donde el equipo evalúe el resultado antes de entregarlo. & \textbf{Casi nunca} & Falta el momento en que el trabajo deja de ser reparto y se vuelve interaprendizaje \citep{bertely2011interaprendizajes}. \\
\midrule
% Escala de frecuencia (comentada, no visible): Siempre / Casi siempre / A veces / Casi nunca / Nunca.
\multicolumn{4}{@{}l}{\cellcolor{unadmRowAlt}\textbf{Dimensión 5. Metacognición y atención a la diversidad}} \\
\midrule
17 & Reviso la retroalimentación recibida y la incorporo en el trabajo siguiente. & \textbf{Casi siempre} & La corrección ocurre, pero producto por producto y sin acumular aprendizaje. \\
\midrule
18 & Llevo un registro de mis avances y dificultades a lo largo del periodo. & \textbf{Casi nunca} & Sin registro no hay patrón visible y cada error se enfrenta como si fuera el primero. \\
\midrule
19 & Reconozco la diversidad cultural y lingüística del grupo como un recurso de aprendizaje. & \textbf{Casi siempre} & El reconocimiento existe en el plano del principio y orienta la disposición general. \\
\midrule
20 & Ajusto el lenguaje y los ejemplos que uso para que sean accesibles a personas con trayectorias distintas. & \textbf{A veces} & La brecha entre reconocer la diversidad y actuar en consecuencia se juega en esta práctica concreta \citep{aguadoSleeter2021educacion}. \\
\end{longtable}
\endgroup
\end{landscape}
\clearpage
\pagestyle{fancy}

\subsection{Perfil que arroja el instrumento}

Leído en conjunto, el registro muestra diez respuestas de frecuencia alta
(``Siempre'' o ``Casi siempre'') y diez de frecuencia intermedia o baja
(``A veces'' o ``Casi nunca''); no se marcó ``Nunca''. En concreto, las frecuencias
altas aparecen en los ítems 1, 2, 5, 6, 9, 10, 13, 14, 17 y 19, mientras que las
respuestas ``A veces'' o ``Casi nunca'' corresponden a los ítems 3, 4, 7, 8, 11,
12, 15, 16, 18 y 20. Por tanto, la tabla permite identificar oportunidades de
mejora en la interacción, la revisión conjunta y el ajuste del lenguaje, pero no
permite atribuir esas diferencias únicamente a la dependencia de otras personas.
Mi interpretación es que el compromiso prioritario no consiste en aumentar la
disciplina general, sino en convertir la autogestión ya declarada en reciprocidad
observable.

Ese diagnóstico tiene una implicación práctica inmediata. Si el vacío está en la
relación y no en el esfuerzo, entonces los compromisos de mejora deben formularse
como conductas observables dentro de la interacción, y no como propósitos de
dedicación. Un compromiso del tipo ``estudiaré más'' sería inútil aquí; uno del tipo
``citaré explícitamente la aportación de otra persona'' es verificable y ataca el
punto débil identificado.

\subsection{Ficha de recuperación y compromisos}

Con esa base se elabora la Ficha de recuperación, que se reproduce a continuación
tal como fue entregada. La ficha recupera los hallazgos por dimensión, sintetiza la
autopercepción y establece dos compromisos: uno personal, orientado a convertir la
participación expositiva en conversacional, y uno colaborativo, orientado a instalar
la revisión conjunta del producto de equipo. Ambos se acompañan de un criterio de
verificación, de modo que en la semana 9 sea posible constatar avances y dificultades
y no solo recordar intenciones.

\begin{fichabox}[title={Ficha de recuperación de mi Brújula de Interaprendizaje}]
\hfill\fichaCopyButton{ficha-brujula-Actividad-4.txt}\\[-0.5em]
\phantomsection\label{ficha-brujula}
\fichameta{Brújula de Interaprendizaje (cuestionario de autopercepción, escala de frecuencia)}{Martin Jonathan de la Cruz}{\today}{Semana 4, aula virtual UnADM}

\textbf{Hallazgos por dimensión.}\par
\medskip
\textit{Autogestión.} Es la dimensión más firme: planifico a partir del calendario y
sostengo un ritmo constante. La debilidad no está en la planeación sino en la
holgura, pues al complicarse una entrega desaparece el margen de reserva y el repaso
se sacrifica antes que el producto.

\textit{Comunicación dialógica.} Participo con regularidad y fundamento lo que
escribo, pero de manera expositiva: aporto mi posición más de lo que pregunto por la
ajena. Leo con atención las aportaciones de otras personas, aunque no siempre las
retomo de forma explícita, que es donde el diálogo se vuelve visible.

\textit{Recursos digitales y aprendizaje ubicuo.} Manejo con soltura las herramientas
y aprovecho tiempos y lugares diversos. El punto débil es la dispersión: acumulo
materiales más rápido de lo que los integro, y la disponibilidad pesa más que el
criterio de selección.

\textit{Colaboración.} Cumplo los acuerdos y entrego lo que me corresponde, pero mi
aportación se concentra en la parte asignada. Reviso poco el producto completo y
participo poco en el cierre colectivo, justo donde el trabajo deja de ser reparto de
tareas para volverse construcción conjunta.

\textit{Metacognición y diversidad.} Reviso la retroalimentación y corrijo, aunque
producto por producto y sin un registro que permita ver el patrón. Reconozco la
diversidad cultural y lingüística como valor, pero rara vez ajusto el lenguaje o los
ejemplos para que el intercambio sea efectivamente accesible.

\medskip
\textbf{Síntesis de la autopercepción.} La brújula apunta a un perfil autogestivo y
constante, con una interacción más individual que dialógica. Soy sólido en lo que
depende de mí y más débil en lo que depende de la relación con otras personas. Lo que
falta no es disciplina, sino convertirla en reciprocidad.

\medskip
\textbf{Compromiso personal.} Antes de publicar cualquier aportación en foro o wiki,
leeré al menos dos aportaciones previas y citaré de forma explícita una de ellas para
continuarla, matizarla o preguntarle algo; además llevaré un registro semanal breve
de lo que aprendí de otra persona. \textit{Verificación:} cada aportación mía debe
contener una referencia nominal a la aportación de alguien más, y el registro debe
tener una entrada por semana, de la semana 5 a la 10.

\medskip
\textbf{Compromiso colaborativo.} En los trabajos de equipo propondré y sostendré una
revisión conjunta del producto completo antes de la entrega, con un turno de palabra
por integrante, y me ofreceré a redactar la síntesis final integrando lo que aportó
cada quien sin sustituir sus palabras, cuidando que el lenguaje resulte comprensible
para personas con trayectorias y lenguas distintas. \textit{Verificación:} constancia
de una sesión o hilo de revisión conjunta documentado antes de cada entrega en
equipo, en las semanas 6 y 8.

\medskip
\textbf{Declaración de uso de inteligencia artificial.} Se utilizó una herramienta de
inteligencia artificial generativa para organizar ideas y revisar la redacción de
esta ficha. Las respuestas al instrumento, los hallazgos y los compromisos
corresponden a mi propia experiencia como estudiante en entornos virtuales.
\end{fichabox}

\clearpage
\section{Conclusión}

El ejercicio confirma que ser estudiante en un entorno virtual es una condición
aprendida y, por lo mismo, corregible. El instrumento fue respondido con criterios
claros: cada afirmación se contrastó con conductas efectivamente realizadas durante
el periodo y no con una imagen deseable de mí mismo, y cada frecuencia marcada se
interpretó desde el marco conceptual de la semana en lugar de leerse como un valor
aislado. Esa disciplina de lectura es lo que convierte una autopercepción en un
diagnóstico utilizable.

A mi juicio, el hallazgo más valioso no es ninguna respuesta en particular, sino su
distribución: todas las fortalezas caen del lado de lo que depende de mí y todas las
debilidades del lado de lo que depende de la relación con otras personas. Ese patrón
explica por qué la autogestión, siendo una virtud, puede volverse un límite cuando
sustituye al diálogo; y explica también por qué los compromisos se formularon como
conductas verificables dentro de la interacción y no como propósitos de esfuerzo.

La implicación profesional es directa. El ejercicio del Derecho rara vez es
solitario: litigar, asesorar o redactar en equipo exige integrar versiones distintas
de un mismo problema y comunicarse con personas cuya lengua, formación y contexto
difieren de los propios. Un profesional que solo sabe cumplir su parte produce
documentos correctos y decisiones parciales. Por eso los dos compromisos asumidos
--retomar explícitamente lo que otras personas aportan y sostener la revisión
conjunta antes de cerrar-- no son hábitos escolares, sino ensayos anticipados de la
forma de trabajo que la práctica jurídica va a exigir.\footnote{Para redactar y
organizar este documento se utilizó una herramienta de inteligencia artificial con el
propósito de ordenar ideas y revisar la redacción; su uso no sustituyó la comprensión
de las lecturas, la aplicación del instrumento ni la elaboración de la ficha y los
compromisos, que corresponden al autor.}

\clearpage
\nocite{bautista2016didactica,borges2007estudiante,zambranoPilay2020tecnologias,aguadoSleeter2021educacion,conocer2025ec0554,bertely2011interaprendizajes}
\bibliography{interaprendizaje-en-ambientes-virtuales}

\end{document}
